\chapter{Dolne twierdzenie Löwenheima--Skolema i domknięcia Skolema}
\label{ch:downward-ls}

\section{Dlaczego nieskończony model może mieć mniejszy model elementarny?}
W rozdziale~\ref{ch:model-morphisms} zdefiniowaliśmy podstrukturę
elementarną. Przypomnijmy: zapis $\mathcal N\preccurlyeq\mathcal M$
oznacza, że $N\subseteq M$, interpretacje symboli w $\mathcal N$
są ograniczeniami interpretacji z $\mathcal M$, a ponadto
dla każdej formuły $\varphi(\vec x)$ języka i każdej
krotki $\vec a\in N^{|\vec x|}$ zachodzi
\[
\mathcal N\models\varphi(\vec a)
\quad\Longleftrightarrow\quad
\mathcal M\models\varphi(\vec a).
\]
Nie wystarcza samo zachowanie działań i relacji: musimy
zachować także prawdziwość zdań egzystencjalnych z parametrami
już należącymi do podstruktury.

Wybór \emph{jednego} świadka dla każdego żądania istnienia
pozwoli nam zbudować stosunkowo małą podstrukturę elementarną.
W dowodach dotyczących dowolnych mocy korzystamy z aksjomatu
wyboru \emph{w metateorii}; nie wprowadzamy go jako aksjomatu
badanej struktury. Aby dokonywać wyborów spójnie, najpierw
dobrze porządkujemy całe uniwersum modelu.

\section{Liczba formuł i wstępne oszacowanie mocy}
\begin{definicja}[Równoliczność i przeliczalność]
\label{def:ls-cardinality}
Dla zbiorów $X$ i $Y$ zapis $|X|=|Y|$ oznacza istnienie
bijekcji $X\to Y$. Moc zbioru $\N$ oznaczamy $\aleph_0$.
Zbiór $X$ nazywamy \emph{przeliczalnie nieskończonym},
gdy $|X|=\aleph_0$, a \emph{co najwyżej przeliczalnym},
gdy jest skończony albo przeliczalnie nieskończony.
Symbole $\kappa$ i $\lambda$ oznaczają niżej moce
nieskończone, jeśli nie określono inaczej.
\end{definicja}

\begin{uwaga}[Arytmetyka mocy a aksjomat wyboru]
Dla każdej mocy nieskończonej $\kappa$, przy aksjomacie
wyboru, zachodzi $\kappa^n=\kappa$ dla $n\ge1$
oraz $\aleph_0\cdot\kappa=\kappa$.
Wykorzystujemy te równości w metateorii: nie są to
równości zwykłych liczb naturalnych. Liczba symboli
pozalogicznych języka $\mathcal L$ będzie oznaczana
przez $|\mathcal L|$.
\end{uwaga}

\begin{lemat}[Liczba formuł z parametrami]
\label{lem:ls-formula-count}
Niech $\mathcal L$ będzie językiem pierwszego rzędu,
a $X$ dowolnym zbiorem. Połóżmy
$\kappa:=\max(\aleph_0,|\mathcal L|,|X|)$.
Zbiór par $(\varphi,\vec a)$, gdzie
$\varphi(\vec x)$ jest formułą języka $\mathcal L$,
a $\vec a\in X^{|\vec x|}$ jest krotką parametrów
odpowiadającą ustalonym wolnym zmiennym formuły,
ma moc co najwyżej $\kappa$.
\end{lemat}
\begin{proof}
Alfabet języka zawiera nie więcej niż
$\max(\aleph_0,|\mathcal L|)\le\kappa$
symboli: oprócz symboli pozalogicznych
występuje przeliczalny zbiór zmiennych
i skończenie wiele rodzajów symboli logicznych.
Formuła jest skończonym napisem, więc
zbiór wszystkich formuł ma moc co najwyżej
$\sum_{n\in\N}\kappa^n=\kappa$.
Zbiór wszystkich skończonych krotek z $X$
ma moc co najwyżej
$\sum_{n\in\N}|X|^n\le\sum_{n\in\N}\kappa^n=\kappa$.
Rozpatrywane pary tworzą podzbiór iloczynu
tych dwóch zbiorów, którego moc jest
co najwyżej $\kappa\cdot\kappa=\kappa$.
Korespondencja pomiędzy liczbą wyróżnionych
zmiennych a długością krotki może jedynie
zmniejszyć ten zbiór.
\end{proof}

\section{Zewnętrzne funkcje wyboru świadków}
\begin{definicja}[Funkcje Skolema dla ustalonej struktury]
\label{def:ls-skolem-functions}
Niech $\mathcal M$ będzie strukturą języka
$\mathcal L$ o niepustym uniwersum $M$.
Dla każdej formuły $\theta(y,\vec x)$, w której
$y$ jest wyróżnioną zmienną wolną,
a $\vec x=(x_1,\ldots,x_k)$ pozostałą
ustaloną listą zmiennych wolnych, wybierzmy funkcję
$F_\theta:M^k\to M$ taką, że dla każdego
$\vec a\in M^k$:
\[
\mathcal M\models\exists y\,\theta(y,\vec a)
\quad\Longrightarrow\quad
\mathcal M\models\theta(F_\theta(\vec a),\vec a).
\]
Gdy świadek nie istnieje, funkcja może przyjąć
dowolną, lecz z góry ustaloną wartość z $M$.
Tak dobrane funkcje nazywamy
\emph{zewnętrznymi funkcjami Skolema}
dla struktury $\mathcal M$.
Nie dodajemy ich symboli do języka $\mathcal L$.
\end{definicja}

\begin{uwaga}[Dlaczego równoczesny wybór jest możliwy?]
\label{rem:ls-choice}
Ustalmy dobry porządek $\triangleleft$ na $M$
i jeden element $m_*\in M$.
Dla danych $\theta,\vec a$ wybieramy jako
$F_\theta(\vec a)$ najmniejszy w tym porządku
element spełniający $\theta(y,\vec a)$,
o ile taki istnieje; w przeciwnym razie
wybieramy $m_*$. W ten sposób określamy całą
rodzinę funkcji, bez osobnych nieuzasadnionych
wyborów na kolejnych etapach.
Funkcje te nie muszą być ani obliczalne,
ani definiowalne przez formuły języka $\mathcal L$.
\end{uwaga}

\begin{definicja}[Domknięcie Skolema]
\label{def:ls-hull}
Niech $\mathcal M$ będzie strukturą o niepustym
uniwersum $M$, niech $A\subseteq M$, a
$m_*\in M$ i funkcje $F_\theta$ niech będą ustalone
jak w poprzedniej definicji i uwadze.
Definiujemy indukcyjnie rosnący ciąg zbiorów:
\[
\begin{aligned}
X_0&:=A\cup\{m_*\}\cup
       \{c^{\mathcal M}:c\text{ jest symbolem stałej }\mathcal L\},\\
X_{n+1}&:=X_n\cup
\{F_\theta(\vec a):
  \theta(y,\vec x)\text{ jest formułą }\mathcal L,\
  \vec a\in X_n^{|\vec x|}\}.
\end{aligned}
\]
Zbiór $\operatorname{Hull}_{\mathcal M}(A):=
\bigcup_{n\in\N}X_n$ nazywamy
\emph{domknięciem Skolema zbioru $A$}
względem ustalonych funkcji.
\end{definicja}

\begin{lemat}[Domknięcie Skolema jest elementarne]
\label{lem:ls-hull-elementary}
Dla struktury $\mathcal M$, zbioru $A\subseteq M$
i funkcji z Definicji~\ref{def:ls-hull},
zbiór $N:=\operatorname{Hull}_{\mathcal M}(A)$
jest uniwersum podstruktury
$\mathcal N\preccurlyeq\mathcal M$ zawierającej $A$.
\end{lemat}
\begin{proof}
Zbiór $N$ jest niepusty, bo zawiera $m_*$,
a wszystkie interpretacje stałych języka
należą do $X_0\subseteq N$.
Uzasadnijmy zamknięcie względem dowolnego
$k$-argumentowego symbolu funkcyjnego $f$.
Jeśli $\vec a=(a_1,\ldots,a_k)\in N^k$,
to dla każdej współrzędnej istnieje indeks $n_i$
z $a_i\in X_{n_i}$. Wybieramy
$n:=\max_i n_i$, a dla $k=0$ wystarcza
wcześniejsze uwzględnienie symboli stałych.
Wówczas $\vec a\in X_n^k$.
Formuła $\theta(y,\vec x):=(y=f(\vec x))$
ma w $\mathcal M$ dokładnie jednego świadka,
zatem $F_\theta(\vec a)=f^{\mathcal M}(\vec a)$.
Definicja $X_{n+1}$ pokazuje, że
$f^{\mathcal M}(\vec a)\in N$.
Możemy więc ograniczyć interpretacje symboli
z $\mathcal M$ do $N$ i otrzymać podstrukturę
$\mathcal N$.

Pozostaje sprawdzić elementarność.
Niech $\eta(y,\vec x)$ będzie dowolną
formułą i $\vec a\in N^{|\vec x|}$.
Jeśli $\mathcal M\models\exists y\,\eta(y,\vec a)$,
dobieramy $n$ tak, by $\vec a\in X_n^{|\vec x|}$.
Z definicji $F_\eta$ mamy
$\mathcal M\models
\eta(F_\eta(\vec a),\vec a)$,
a z konstrukcji $F_\eta(\vec a)\in X_{n+1}\subseteq N$.
Każda prawdziwa w $\mathcal M$ formuła egzystencjalna
z parametrami z $N$ ma więc w $N$ świadka.
Kryterium Tarskiego--Vaughta,
Twierdzenie~\ref{thm:model-tarski-vaught},
daje $\mathcal N\preccurlyeq\mathcal M$.
\end{proof}

\section{Dolne twierdzenie Löwenheima--Skolema}
\begin{twierdzenie}[Dolne twierdzenie Löwenheima--Skolema]
\label{thm:downward-ls}
Niech $\mathcal L$ będzie językiem pierwszego
rzędu, $\mathcal M$ strukturą tego języka
o niepustym uniwersum $M$, a $A\subseteq M$.
Połóżmy $\kappa:=\max(\aleph_0,|\mathcal L|,|A|)$.
Wówczas istnieje podstruktura elementarna
$\mathcal N\preccurlyeq\mathcal M$, dla której
$A\subseteq N$ oraz $|N|\le\kappa$.

Jeżeli ponadto $\mathcal M$ jest nieskończona,
to dla każdej mocy nieskończonej $\lambda$
spełniającej $\kappa\le\lambda\le |M|$
istnieje taka $\mathcal N$ o mocy
dokładnie $|N|=\lambda$.
\end{twierdzenie}
\begin{proof}
Wybierzmy dobry porządek uniwersum $M$
i skonstruujmy zewnętrzne funkcje Skolema
z Definicji~\ref{def:ls-skolem-functions}.
Dla $A$ budujemy ciąg $X_n$
z Definicji~\ref{def:ls-hull}.
Wykażemy indukcyjnie $|X_n|\le\kappa$.
Dla $n=0$ mamy
$|X_0|\le |A|+1+|\mathcal L|\le\kappa$,
gdyż interpretacji stałych jest co najwyżej
tyle, ile symboli języka.
Jeśli $|X_n|\le\kappa$, to
z Lematu~\ref{lem:ls-formula-count}
wynika, że istnieje co najwyżej
$\kappa$ par $(\theta,\vec a)$ z parametrami
w $X_n$. Z każdej pary dodajemy
co najwyżej jeden element, dlatego
$|X_{n+1}|\le\kappa+\kappa=\kappa$.
Wobec tego
$|\bigcup_{n\in\N}X_n|
\le\aleph_0\cdot\kappa=\kappa$.
Lemat~\ref{lem:ls-hull-elementary}
gwarantuje elementarność otrzymanej
podstruktury i zawieranie zbioru $A$.

Dla drugiej części wybierzmy zbiór
$A_0\subseteq M$ taki, że
$A\subseteq A_0$ oraz $|A_0|=\lambda$.
Jest to możliwe, ponieważ
$|A|\le\lambda\le|M|$ i korzystamy
z aksjomatu wyboru w metateorii.
Pierwsza część twierdzenia, zastosowana
do $A_0$, daje podstrukturę elementarną
$\mathcal N$, dla której
$A_0\subseteq N$ i
\[
|N|\le\max(\aleph_0,|\mathcal L|,|A_0|)
     =\lambda.
\]
Z drugiej strony $|N|\ge|A_0|=\lambda$,
więc $|N|=\lambda$.
\end{proof}

\begin{wniosek}[Przeliczalny podmodel elementarny]
\label{cor:downward-countable}
Jeżeli $\mathcal L$ jest co najwyżej przeliczalnym
językiem pierwszego rzędu, a $\mathcal M$
nieskończoną strukturą tego języka,
to istnieje przeliczalnie nieskończona
podstruktura elementarna
$\mathcal N\preccurlyeq\mathcal M$.
Możemy wymagać dodatkowo, aby $N$
zawierał dowolny wskazany
co najwyżej przeliczalny zbiór $A\subseteq M$.
\end{wniosek}
\begin{proof}
W Twierdzeniu~\ref{thm:downward-ls}
wybieramy $\lambda=\aleph_0$.
Wówczas $\kappa=\aleph_0$, a
$\aleph_0\le|M|$ na mocy nieskończoności
$\mathcal M$ i aksjomatu wyboru.
Otrzymujemy $A\subseteq N$ oraz $|N|=\aleph_0$.
\end{proof}

\section{Przykłady, kontrprzykłady i ograniczenia}
\begin{przyklad}[Przeliczalna podstruktura elementarna liczb rzeczywistych]
\label{ex:ls-real-closed}
Rozważmy strukturę uporządkowanego ciała
$\mathcal R=(\R,0,1,+,\cdot,<)$.
Jej język jest skończony. Wniosek~\ref{cor:downward-countable}
daje więc przeliczalnie nieskończoną
$\mathcal K\preccurlyeq\mathcal R$.
Podstruktura w rozpatrywanej sygnaturze
zawiera $0,1$ i jest zamknięta
względem dodawania i mnożenia.
Elementarność daje znacznie więcej.
Ponieważ $\mathcal R\models\exists x\,(x\cdot x=1+1)$,
to również $\mathcal K$ zawiera pierwiastek
z liczby $2$ (a także jego element przeciwny).
Ponadto dla każdego $a\in K$
formuła $\exists y\,(a+y=0)$ ma świadka
w $\mathcal R$, zatem ma go w $\mathcal K$;
dla $a\ne0$ to samo dotyczy
$\exists y\,(a\cdot y=1)$.
Zatem $K$ jest ciałem, mimo że
w sygnaturze nie zapisaliśmy osobnych symboli
odejmowania i dzielenia.
Jest to przeliczalna podstruktura elementarna
nieprzeliczalnej struktury: nie oznacza to
przeliczalności całego $\R$.
\end{przyklad}

\begin{przyklad}[Dlaczego $\Q$ nie jest podstrukturą elementarną $\R$?]
\label{ex:ls-q-not-elementary}
W języku $(0,1,+,\cdot,<)$ struktura
$\mathcal Q=(\Q,0,1,+,\cdot,<)$
jest podstrukturą $\mathcal R$,
ale nie jest elementarna.
Zdanie $\exists x\,(x\cdot x=1+1)$
jest prawdziwe w $\mathcal R$,
a fałszywe w $\mathcal Q$.
Gdyby $x=p/q$ dla względnie pierwszych
$p\in\Z$, $q\in\Z\setminus\{0\}$
spełniało $x^2=2$, mielibyśmy
$p^2=2q^2$. Parzystość $p^2$
wymusza parzystość $p$, a po podstawieniu
$p=2r$ otrzymalibyśmy $q^2=2r^2$,
czyli również parzystość $q$.
To przeczy względnej pierwszości $p,q$.
Zamknięcie na działania nie wystarcza więc
do zachowania formuł egzystencjalnych.
\end{przyklad}

\begin{uwaga}[Dwie istotne granice twierdzenia]
\label{rem:ls-limitations}
Jeżeli język ma nieprzeliczalnie wiele
symboli stałych, interpretowanych jako różne
elementy, każda podstruktura musi je wszystkie
zawierać. Nie można wtedy oczekiwać
przeliczalnego podmodelu elementarnego.
Ponadto opisana konstrukcja nie jest algorytmem
wyznaczającym mały model z dowolnego
opisu modelu dużego: dobór
$F_\theta(\vec a)$ wymaga wiedzy
o prawdziwości odpowiedniej formuły
egzystencjalnej w $\mathcal M$, której
nie musimy umieć rozstrzygać.
\end{uwaga}
